\section{Lean and ConciseIPRA formalization plan}
\label{sec:lean}

\subsection{Repository interfaces}

The intended target is the checked embedded calculus used by ProveIt's
\code{IPRADerivations} library.  A finished theorem has the shape
\[
 ([]:\operatorname{List}(\code{Fm PTm 0}))
 \ \vdash[\code{PRAAx}]\ \varphi.
\]
The checked arithmetic modules already provide:

\begin{itemize}
\item primitive-recursive definitions and defining equations;
\item semiring laws, cancellation, order, and decidable rational-code comparisons;
\item finite sums, finite products, powers, and factorials;
\item pairing, simultaneous recursion, division with remainder, and bounded search;
\item product invariance under an explicitly invertible permutation;
\item a tactic layer for polynomial identities over encoded integers.
\end{itemize}

Once an ordinary iPRA derivation exists, \code{pra_iff_native} in
\code{ConciseIPRA.NativeBiInterpretation} transports it to the native scalar core.  The bridge does not create the derivation; it is the final checked translation.

\subsection{Host-level organization types}

The companion file \code{lean/BetaIPRAOutline.lean} uses ordinary Lean structures to organize obligations.  These structures are metatheoretic bookkeeping, not new sorts added to iPRA.

\begin{lstlisting}[language={}]
structure RatCode where
  num : Int
  den : Nat
  den_pos : 0 < den

structure CQRat where
  re : RatCode
  im : RatCode

structure RegularName where
  approx : Nat -> RatCode
  regular : forall p q,
    absQ (approx p - approx q) <= dyadic p + dyadic q

structure BetaDomain (a b : CName) where
  A B R : Nat
  A_pos : 0 < A
  B_pos : 0 < B
  re_a_margin : Re a >= 1 / A
  re_b_margin : Re b >= 1 / B
  norm_a : norm1 a <= R
  norm_b : norm1 b <= R
\end{lstlisting}

In the actual object theory, a \code{RegularName} is not quantified over.  A particular \code{approx} implementation is compiled to a primitive-recursive symbol, and its regularity theorem is an assumption or previously derived theorem about that symbol.

\subsection{The concrete beta evaluator}

For domain data $D=(A,B,R)$ and precision $p$, define:
\begin{align*}
 L(D,p)&=p+A+B+10,\\
 h(D,p)&=AB\,L(D,p),\\
 \varepsilon(D,p)&=2^{-h(D,p)},\\
 C(D,p)&=2(R+1)\varepsilon(D,p)^{-2R-3},\\
 m(D,p)&=2^{p+7}(1+\lceil C(D,p)\rceil).
\end{align*}
The evaluator is the left sum
\[
 \operatorname{betaApprox}(D,a,b,p)
 =\frac{1-2\varepsilon}{m}
 \sum_{j=0}^{m-1}
 f_{a,b}\!\left(
 \varepsilon+\frac{j(1-2\varepsilon)}{m}\right),
\]
with each elementary-function call made at the precision allocated by the product-error lemma.  All divisions are rational-code divisions with visibly positive denominators.

The first Lean milestones are:

\begin{enumerate}
\item prove that every sample point lies in
$[\varepsilon,1-\varepsilon]$;
\item prove the endpoint bounds
\eqref{eq:lefttail}--\eqref{eq:righttail};
\item prove the derivative majorant $C(D,p)$;
\item prove that the Riemann error is below $2^{-p-7}$;
\item prove regularity and independence from enlarged valid domain data.
\end{enumerate}

The internal formula for regularity is a cleared integer inequality.  It contains no reference to an abstract integral.

\subsection{Module dependency order}

\begin{longtable}{@{}p{0.21\textwidth}p{0.46\textwidth}p{0.25\textwidth}@{}}
\toprule
Planned module & Main declarations & iPRA invariant\\
\midrule\endhead
\code{Beta/RatCode} &
Signed fractions, dyadics, complex pairs, absolute-value comparisons, error propagation &
Integer cross products; positivity of denominators\\

\code{Beta/Elementary} &
Certified \code{exp}, \code{log}, positive powers, sine, cosine, square root and positive rational powers &
Finite Taylor/geometric sums and factorial tails\\

\code{Beta/Names} &
Regular-name equality, arithmetic, guarded inverse, substitution of fixed names &
Precision fixed before induction; rational inequalities only\\

\code{FiniteIntegral} &
Riemann sums, common refinement, finite fundamental theorem, integration by parts, substitution &
Telescoping sums with a quadratic remainder\\

\code{Beta/Complete} &
\code{betaApprox}, endpoint and mesh certificates, symmetry, positivity &
Mesh-index induction and list reversal\\

\code{Beta/Recurrence} &
Pascal, balanced integration by parts, coordinate shifts, factorial values &
Finite product concatenation and endpoint terms\\

\code{Beta/GammaBridge} &
Effective limit \eqref{eq:beta-gamma-limit}, exact Euler ratio, beta--gamma &
Finite logarithm inequalities and Gamma approximant error\\

\code{Beta/Transforms} &
Trigonometric, infinite-interval, power, fractional-linear, central radical forms &
Substitution telescope plus endpoint budgets\\

\code{Beta/SeriesProduct} &
Binomial series, Euler product, product tails &
Finite Taylor polynomial; product telescope; $\sum k^{-2}$ bound\\

\code{Beta/Derivatives} &
Logarithmic moments and digamma/polygamma formulas &
Finite parameter Taylor polynomial\\

\code{Multiplication} &
General multiplication and Gosper odd/even products &
Permutation of residues; carry-count product\\

\code{Beta/Incomplete} &
Incomplete evaluator, recurrences, polynomial values, hypergeometric name, binomial CDF &
Finite interval split; binomial theorem; polynomial telescoping\\

\code{ContinuedFraction} &
Continuants, determinant identity, positive remainder &
Second-order recursion and sign invariant\\

\code{Beta/Multi} &
List Gamma quotient, shifts, Pascal, fixed-dimensional simplex &
List sum/product recursion; iterated one-dimensional beta reduction\\

\code{Beta/Asymptotic} &
Stirling combination, fixed-parameter limit &
Addition of three explicit Gamma remainders\\

\code{Beta/ObjectTheory} &
Reified \code{Fm PTm} statements and \code{Prf PRAAx} terms &
Only \code{PRA.qfInd} on encoded finite invariants\\
\bottomrule
\end{longtable}

\subsection{Representative theorem statements in Lean}

The metatheoretic analysis layer can expose the following shapes.  The exact names depend on the checked-out ProveIt revision.

\begin{lstlisting}[language={}]
theorem beta_symm
    (D : BetaDomain a b) :
    NameEq (betaName D a b) (betaName D.swap b a) := by
  -- reverse the common Riemann mesh; discharge tails
  ...

theorem beta_pascal
    (D : BetaDomain a b) :
    NameEq (betaName D a b)
      (betaName (D.shiftLeft) (a+1) b
       + betaName (D.shiftRight) a (b+1)) := by
  -- pointwise identity t + (1-t) = 1
  ...

theorem beta_shift_left
    (D : BetaDomain a b) :
    NameEq (betaName D.shiftLeft (a+1) b)
      (a / (a+b)) * betaName D a b := by
  -- Pascal plus finite integration by parts
  ...

theorem beta_gamma
    (D : BetaDomain a b)
    (G : EffectiveGammaContract) :
    NameEq (betaName D a b)
      (G.gamma a * G.gamma b / G.gamma (a+b)) := by
  -- effective N^a B(a,b+N) limit and exact Euler ratio
  ...

theorem incomplete_beta_nat
    (m n : Nat) (hm : 0 < m) (hn : 0 < n) :
    PolyEq (incBetaPoly m n)
      (sum n (fun j =>
        (-1)^j * choose (n-1) j * X^(m+j) / (m+j))) := by
  -- finite binomial theorem
  ...
\end{lstlisting}

The delivered outline keeps the analytic holes visible as \code{sorry}; it is not represented as a completed Lean proof.

\subsection{Object-theory formulas}

A name equality must be turned into an arithmetical sentence.  Let
$\mathsf{BA}(D,a,b,p)$ be the rational-code output of the evaluator and let
$\operatorname{close}(u,v,p)$ be the Boolean relation
$|u-v|\le2^{-p}$ after positive denominators are cleared.  The symmetry target is
\[
 \forall D,a,b,p\;
 \operatorname{Dom}(D,a,b)\to
 \operatorname{close}\bigl(
 \mathsf{BA}(D,a,b,p+c),
 \mathsf{BA}(D^\mathrm{op},b,a,p+c),p\bigr).
\]
Here $D,a,b$ are finite tuples of natural-number codes, and
$\operatorname{Dom}$ is a primitive-recursive Boolean conjunction.

The Pascal target is
\[
 \operatorname{close}\bigl(
 \mathsf{BA}(a,b,p+c),
 \mathsf{BA}(a+1,b,p+c)+
 \mathsf{BA}(a,b+1,p+c),p\bigr).
\]
At a common mesh, the summand identity is exact.  The proof object consists of:
\begin{enumerate}
\item the derivation of the summand identity by the ring tactic;
\item finite-sum congruence;
\item the endpoint and mesh inequalities;
\item integer triangle inequalities.
\end{enumerate}

For the beta--gamma identity, division is avoided:
\[
 \operatorname{close}\bigl(
 \mathsf{BA}(a,b,p+c)\cdot\mathsf{GA}(a+b,p+c),
 \mathsf{GA}(a,p+c)\cdot\mathsf{GA}(b,p+c),p\bigr).
\]
Only after this theorem is available does a guarded-inverse lemma derive the quotient form.

\subsection{Quantifier-free induction invariants}

The important invariants are concrete.

\begin{description}
\item[Mesh reversal.]  The sum of the first $r$ beta samples for $(a,b)$ equals the sum of the last $r$ reversed samples for $(b,a)$, up to a tracked rational evaluation error.
\item[Pascal.]  The difference between the accumulated left side and the two accumulated right sides is exactly zero for the first $r$ samples.
\item[Integration by parts.]  The accumulated discrete derivative equals the endpoint difference plus a remainder bounded by $rCh^2$.
\item[Rising products.]  After $r$ shifts,
$(a+b)_r\B(a+r,b)=(a)_r\B(a,b)$ at the chosen finite precision.
\item[Odd Gosper permutation.]  The multiplicity counter for residue $j$ among $2i\bmod(2n+1)$ is $1$ after all indices have been processed; the carry product contains exactly the odd residues.
\item[Incomplete polynomial.]  The derivative of the partial polynomial through term $r$ equals the first $r$ terms of the finite binomial expansion.
\item[Binomial CDF.]  The derivative of the first $r$ summands telescopes to two boundary terms; after $r=k+1$ only the final boundary remains.
\end{description}

Each invariant is an equality or inequality among primitive-recursive natural-number terms.  No induction formula contains the unbounded precision quantifier used in name equality.

\subsection{Use of the checked ProveIt arithmetic library}

The most direct imports for the eventual implementation are expected to include
\begin{lstlisting}[language={}]
import ConciseIPRA.NativeBiInterpretation
import IPRADerivations.Core.Basic
import IPRADerivations.Core.Definitions
import IPRADerivations.Core.Sums
import IPRADerivations.Core.Order
import IPRADerivations.Core.Powers
import IPRADerivations.Core.Division
import IPRADerivations.Core.Search
import IPRADerivations.Core.Permutation
\end{lstlisting}
Exact module paths must be adjusted to the current workspace.  The Gamma packet's host-level \code{GammaIPRA.lean} is useful as an interface sketch, but its analytic certificate structures are uninstantiated and should not be imported as if they were proved theorems.

\subsection{Definition of done}

A complete formalization should satisfy all of the following.

\begin{enumerate}
\item Every theorem exported as a beta identity has an explicit term of type
$[]\vdash[\code{PRAAx}]\varphi$.
\item A \texttt{\#print axioms} audit shows no \code{sorryAx}; standard Lean metatheory axioms are listed separately.
\item Every analytic evaluator is a fixed primitive-recursive symbol with defining equations.
\item Every use of inverse includes a proved positive margin.
\item Every improper integral, infinite sum, product, derivative, or asymptotic statement has a closed cutoff function.
\item The final native theorem is obtained only by applying the checked
\code{pra_iff_native} bridge to the already constructed ordinary derivation.
\end{enumerate}

The current archive meets the mathematical-specification and outline requirements, not this final kernel-checked definition of done.
